% calcu3.tex — Preface (source: calcu3.pdf).
%
% Four preface sections (first/second/third/fourth edition), each opening with
% an all-caps heading and closing with a signed location/date. Reproduced as
% flowing prose with a fresh page before each edition, as in the source.

\thispagestyle{empty}
\vspace*{3cm}
\begin{center}
  {\scshape\Huge Preface}
\end{center}
\vspace{3cm}

\epigraph{%
  I hold every man a debtor\\
  to his profession,\\
  from the which as men of course\\
  doe seeke to receive countenance and profit,\\
  so ought they of duty to endeavour\\
  themselves by way of amends,\\
  to be a help and ornament thereunto.%
}{Francis Bacon}

\section*{\mdseries PREFACE TO THE FIRST EDITION}

Every aspect of this book was influenced by the desire to present calculus not
merely as a prelude to but as the first real encounter with mathematics. Since
the foundations of analysis provided the arena in which modern modes of
mathematical thinking developed, calculus ought to be the place in which to
expect, rather than avoid, the strengthening of insight with logic. In addition
to developing the students' intuition about the beautiful concepts of analysis,
it is surely equally important to persuade them that precision and rigor are
neither deterrents to intuition, nor ends in themselves, but the natural medium
in which to formulate and think about mathematical questions.

This goal implies a view of mathematics which, in a sense, the entire book
attempts to defend. No matter how well particular topics may be developed,
the goals of this book will be realized only if it succeeds as a whole. For this
reason, it would be of little value merely to list the topics covered, or to
mention pedagogical practices and other innovations. Even the cursory glance
customarily bestowed on new calculus texts will probably tell more than any
such extended advertisement, and teachers with strong feelings about particular
aspects of calculus will know just where to look to see if this book fulfills
their requirements.

A few features do require explicit comment, however. Of the twenty-nine
chapters in the book, two (starred) chapters are optional, and the three
chapters comprising Part V have been included only for the benefit of those
students who might want to examine on their own a construction of the real
numbers. Moreover, the appendices to Chapters 3 and 11 also contain optional
material.

The order of the remaining chapters is intentionally quite inflexible, since
the purpose of the book is to present calculus as the evolution of one idea,
not as a collection of ``topics.'' Since the most exciting concepts of calculus
do not appear until Part III, it should be pointed out that Parts I and II
will probably require less time than their length suggests---although the
entire book covers a one-year course, the chapters are not meant to be covered
at any uniform rate. A rather natural dividing point does occur between
Parts II and III, so it is possible to reach differentiation and integration
even more quickly by treating Part II very briefly, perhaps returning later for
a more detailed treatment. This arrangement corresponds to the traditional
organization of most calculus courses, but I feel that it will only diminish
the value of the book for students who have seen a small amount of calculus
previously, and for bright students with a reasonable background.

The problems have been designed with this particular audience in mind. They
range from straightforward, but not overly simple, exercises which develop
basic techniques and test understanding of concepts, to problems of
considerable difficulty and, I hope, of comparable interest. There are about
625 problems in all. Those which emphasize manipulations usually contain many
examples, numbered with small Roman numerals, while small letters are used to
label interrelated parts in other problems. Some indication of relative
difficulty is provided by a system of starring and double starring, but there
are so many criteria for judging difficulty, and so many hints have been
provided, especially for harder problems, that this guide is not completely
reliable. Many problems are so difficult, especially if the hints are not
consulted, that the best of students will probably have to attempt only those
which especially interest them; from the less difficult problems it should be
easy to select a portion which will keep a good class busy, but not
frustrated. The answer section contains solutions to about half the examples
from an assortment of problems that should provide a good test of technical
competence. A separate answer book contains the solutions of the other parts
of these problems, and of all the other problems as well. Finally, there is a
Suggested Reading list, to which the problems often refer, and a glossary of
symbols.

I am grateful for the opportunity to mention the many people to whom I owe my
thanks. Jane Bjorkgren performed prodigious feats of typing that compensated
for my fitful production of the manuscript. Richard Serkey helped collect the
material which provides historical sidelights in the problems, and Richard
Weiss supplied the answers appearing in the back of the book. I am especially
grateful to my friends Michael Freeman, Jay Goldman, Anthony Phillips, and
Robert Wells for the care with which they read, and the relentlessness with
which they criticized, a preliminary version of the book. Needless to say,
they are not responsible for the deficiencies which remain, especially since I
sometimes rejected suggestions which would have made the book appear suitable
for a larger group of students. I must express my admiration for the editors
and staff of W.~A. Benjamin, Inc., who were always eager to increase the
appeal of the book, while recognizing the audience for which it was intended.

The inadequacies which preliminary editions always involve were gallantly
endured by a rugged group of freshmen in the honors mathematics course at
Brandeis University during the academic year 1965--1966. About half of this
course was devoted to algebra and topology, while the other half covered
calculus, with the preliminary edition as the text. It is almost obligatory in
such circumstances to report that the preliminary version was a gratifying
success. This is always safe---after all, the class is unlikely to rise up in
a body and protest publicly---but the students themselves, it seems to me,
deserve the right to assign credit for the thoroughness with which they
absorbed an impressive amount of mathematics. I am content to hope that some
other students will be able to use the book to such good purpose, and with
such enthusiasm.

\begin{flushright}
  \begin{minipage}{0.4\textwidth}
    \noindent Waltham, Massachusetts\\[0.5em]
    February 1967\\[2em]
    \scshape Michael Spivak
  \end{minipage}
\end{flushright}

\clearpage
\section*{\mdseries PREFACE TO THE SECOND EDITION}

I have often been told that the title of this book should really be something
like ``An Introduction to Analysis,'' because the book is usually used in
courses where the students have already learned the mechanical aspects of
calculus---such courses are standard in Europe, and they are becoming more
common in the United States. After thirteen years it seems too late to change
the title, but other changes, in addition to the correction of numerous
misprints and mistakes, seemed called for.

There are now separate Appendices for many topics that were previously
slighted: polar coordinates, uniform continuity, parameterized curves,
Riemann sums, and the use of integrals for evaluating lengths, volumes and
surface areas. A few topics, like manipulations with power series, have been
discussed more thoroughly in the text, and there are also more problems on
these topics, while other topics, like Newton's method and the trapezoid rule
and Simpson's rule, have been developed in the problems. There are in all
about 160 new problems, many of which are intermediate in difficulty between
the few routine problems at the beginning of each chapter and the more
difficult ones that occur later.

Most of the new problems are the work of Ted Shifrin. Frederick Gordon pointed
out several serious mistakes in the original problems, and supplied some
non-trivial corrections, as well as the neat proof of Theorem 12-2, which took
two Lemmas and two pages in the first edition. Joseph Lipman also told me of
this proof, together with the similar trick for the proof of the last theorem
in the Appendix to Chapter 11, which went unproved in the first edition.
Roy O.~Davies told me the trick for Problem 11-66, which previously was proved
only in Problem 20-8 [21-8 in the third edition], and Marina Ratner suggested
several interesting problems, especially ones on uniform continuity and
infinite series. To all these people go my thanks, and the hope that in the
process of fashioning the new edition their contributions weren't too badly
botched.

\begin{flushright}
  \begin{minipage}{0.4\textwidth}
    \vspace{2em}
    \scshape Michael Spivak
  \end{minipage}
\end{flushright}

\clearpage
\section*{\mdseries PREFACE TO THE THIRD EDITION}

The most significant change in this third edition is the inclusion of a new
(starred) Chapter 17 on planetary motion, in which calculus is employed for a
substantial physics problem.

In preparation for this, the old Appendix to Chapter 4 has been replaced by
three Appendices: the first two cover vectors and conic sections, while polar
coordinates are now deferred until the third Appendix, which also discusses
the polar coordinate equations of the conic sections. Moreover, the Appendix
to Chapter 12 has been extended to treat vector operations on vector-valued
curves.

Another large change is merely a rearrangement of old material: ``The
Cosmopolitan Integral,'' previously a second Appendix to Chapter 13, is now an
Appendix to the chapter on ``Integration in Elementary Terms'' (previously
Chapter 18, now Chapter 19); moreover, those problems from that chapter which
used the material from that Appendix now appear as problems in the newly
placed Appendix.

A few other changes and renumbering of Problems result from corrections, and
elimination of incorrect problems.

I was both startled and somewhat dismayed when I realized that after allowing
13 years to elapse between the first and second editions of the book, I have
allowed another 14 years to elapse before this third edition. During this time
I seem to have accumulated a not-so-short list of corrections, but no longer
have the original communications, and therefore cannot properly thank the
various individuals involved (who by now have probably lost interest anyway).
I have had time to make only a few changes to the Suggested Reading, which
after all these years probably requires a complete revision; this will have to
wait until the next edition, which I hope to make in a more timely fashion.

\begin{flushright}
  \begin{minipage}{0.4\textwidth}
    \vspace{2em}
    \scshape Michael Spivak
  \end{minipage}
\end{flushright}

\clearpage
\section*{\mdseries PREFACE TO THE FOURTH EDITION}

Promises, promises! In the preface to the third edition I noted that it was 13
years between the first and second editions, and then another 14 years before
the third, expressing the hope that the next edition would appear sooner.
Well, here it is another 14 years later before the fourth, and presumably
final, edition.

Although small changes have been made to some material, especially in Chapters
5 and 20, this edition differs mainly in the introduction of additional
problems, a complete update of the Suggested Reading, and the correction of
numerous errors. These have been brought to my attention over the years by,
among others, Nils von Barth; Philip Loewen; Fernando Mejias; Lance Miller,
who provided a long list, particularly for the answer book; and Michael
Maltenfort, who provided an amazingly extensive list of misprints, errors, and
criticisms.

Most of all, however, I am indebted to my friend Ted Shifrin, who has been
using the book for the text in his renowned course at the University of
Georgia for all these years, and who prodded and helped me to finally make
this needed revision. I must also thank the students in his course this last
academic year, who served as guinea pigs for the new edition, resulting, in
particular, in the current proof in Problem 8-20 for the Rising Sun Lemma,
far simpler than Reisz's original proof, or even the proof in [38] of the
Suggested Reading, which itself has now been updated considerably, again with
great help from Ted.

\begin{flushright}
  \begin{minipage}{0.4\textwidth}
    \vspace{2em}
    \scshape Michael Spivak
  \end{minipage}
\end{flushright}

\clearpage
